\def\baselinestretch{1}

\chapter{Control Objectives}\label{ch:control_objectives}

\def\baselinestretch{1.66}

\section{Introduction}\label{sec:control_introduction}

In Chapter~\ref{ch:introduction} we identified overlay and throughput as the two primary drivers of scanner design, followed in Chapter~\ref{ch:model} by the modeling of its multi-body dynamics. This chapter focuses on mapping those industrial requirements into quantitative objectives that any candidate reference trajectory must satisfy. Here we explore in detail how the synchronization error arises between the wafer and the reticle stages, examine the filtered metrics through which the industry evaluates it, and establish the formal statement of the optimization problem.

\section{Synchronization Error}\label{sec:control_synchronization}

During the scanning exposure, the projection optics image the reticle onto the wafer with a demagnification factor of $\alpha = 0.25$. Thus, the image of the reticle pattern is stationary on the wafer surface only if the wafer follows the reticle motion scaled by $\alpha$ at every instant. The degradation of the printed pattern is therefore not due to the individual tracking error of either stage, but their scaled difference. Let $r_w(t)$ and $r_r(t)$ denote the reference positions of the wafer and reticle stages respectively along the scanning direction, and $y_w(t)$, $y_r(t)$ their measured positions. The \textit{synchronization error} is defined as~\cite{Butler2011_PositionControl}
\begin{equation}
    e_{syn}(t) = \big(y_w(t) - r_w(t)\big) - \alpha \big(y_r(t) - r_r(t)\big).
    \label{eq:esyn}
\end{equation}
Naturally, a common-mode mode disturbance that displaces both stages consistently with the projection ratio leaves $e_{syn}$ unaffected. Conversely, any differential motion between the chains --- exacerbated by the lightly damped structural resonance --- is transferred directly to the printed image.

\section{Exposure Metrics: MA and MSD}\label{sec:control_ma_msd}

For an area of photoresist to become exposed, the illumination dose must overcome a threshold. Despite direct exposure during the scanning motion, the photoresist on the wafer does not respond to the instantaneous value of $e_{syn}$ but to its aggregate effect over the time a given point stays below the light source. Therefore, industry practice evaluates the synchronization error through two moving-window filters matched to the exposure time $t_{exp}$ of one point on the wafer~\cite{Butler2011_PositionControl}:

\begin{itemize}
    \item First, the \textbf{moving average}
    \begin{equation}
        \mathrm{MA}(t_k) = \frac{1}{N}\sum_{n=0}^{N-1} e_{syn}(t_{k-n})
        \label{eq:ma}
    \end{equation}
    measures the mean displacement of the image and manifests as \textit{overlay} error or the misalignment of the printed pattern with respect to the previous layer.

    \item Second, the \textbf{moving standard deviation}
    \begin{equation}
        \mathrm{MSD}(t_k) = \sqrt{\frac{1}{N}\sum_{n=0}^{N-1} \Big(e_{syn}(t_{k-n}) - \mathrm{MA}(t_k)\Big)^{2}}
        \label{eq:msd}
    \end{equation}
    measures the smearing of the aerial image around that mean position and manifests as \textit{fading}, a loss of contrast that distorts the printed feature.
\end{itemize}

For the 38\,nm half-pitch process, the specifications adopted~\cite{Al-Rawashdeh2022_Caracterization} are
\begin{equation}
    \mathrm{MA} \in [-1.25,\ +2]\ \mathrm{nm}, \qquad \mathrm{MSD} \leq 7\ \mathrm{nm},
    \label{eq:ma_msd_spec}
\end{equation}
evaluated at every instant during the exposure window~\cite{Butler2011_PositionControl}. These bounds must be met while the stage sustains the required scanning velocity, operating within a settling time budget of under $10\,\mathrm{ms}$ to preserve wafer throughput. On the two-stage machine modelled here they are not a fixed pass/fail but the tight end of a trade curve: Chapter~\ref{ch:experiments} shows that a sufficiently smooth fourth-order trajectory meets them --- both MSD and MA, to the nanometer --- at the cost of a longer scan, so the specification and the throughput budget are the two ends of the optimization this thesis performs.

\section{Problem Statement}\label{sec:control_problem}

The reference trajectories studied here are fourth-order profiles parameterized by the kinematic bounds ($v_{\max}$, $a_{\max}$, $j_{\max}$, $s_{\max}$) on velocity, acceleration, jerk, and snap respectively. These four variables set both the duration of the motion --- and hence throughput --- and the frequency content of the excitation injected into the machine structure, and thereby the synchronization error during exposure.

Conventional trajectory design evaluates candidates using positioning metrics like settling time, peak error, or the MA/MSD bounds of~\eqref{eq:ma_msd_spec}. Such metrics make it possible to predict which dies will be defective, but reveal nothing about the \textit{morphology} of the failures that accumulate across the wafer. Wafer fabs routinely exploit the converse relationship to diagnose which process step caused a particular spatial pattern of failing dies (the \textit{fingerprint})~\cite{Lam2015_Fingerprint}. However, the forward direction remains an open problem as no published methodology predicts, from a given set of trajectory parameters, the defect pattern they will produce, nor uses that prediction as the cost function during trajectory optimization. This thesis addresses precisely that forward direction.

To evaluate the impact of a candidate trajectory on the overall process yield, the resulting wafer defect map is classified by a convolutional neural network trained on the WM811K dataset~\cite{Wu2015_WM811K}. The network estimates a probability $P_c$ for each defect morphology class $c$, and the trajectory cost is defined as the weighted sum of these probabilities:
\begin{equation}
    \mathcal{J} = \sum_{c \in \mathcal{C}} w_c P_c,
    \label{eq:cost}
\end{equation}
where $w_c$ is the economic weight of class $c$~\cite{Kim2017_Market, Kim2019_Productivity, Yang2018_Yield} and $\mathcal{C}$ is the full set of defect-morphology classes. Scanner motion can plausibly produce only a subset of them (namely \textit{none}, \textit{Scratch}, \textit{Edge-Loc}, and \textit{Random}), priced as a tradeable yield loss; the remaining classes (\textit{Center}, \textit{Donut}, \textit{Edge-Ring}, \textit{Loc}, \textit{Near-full}) originate in subsequent steps of the manufacturing process and are assigned a separate, heavier weight that acts as a hard wall rather than a tradeable cost (Section~\ref{sec:methodology_pso}). If they were to be activated by the trajectory alone, it would indicate a modeling error rather than a trajectory effect; this plausible/implausible split is therefore used as a validation criterion for the trajectory-to-defect mapping in Chapter~\ref{ch:experiments}.

The optimization problem is then formulated as
\begin{equation}
    \min_{(v_{\max},\, a_{\max},\, j_{\max},\, s_{\max})} \ \mathcal{J}
    \qquad \text{subject to} \qquad T_{wafer} \leq T_{budget},
    \label{eq:problem}
\end{equation}
where $T_{wafer}$ is the total time to traverse a wafer determined by the trajectory parameters, and $T_{budget}$ is the cycle-time budget corresponding to the target throughput. The central hypothesis of this work is that~\eqref{eq:problem} can be solved using a data-driven cost function of the form~\eqref{eq:cost}, producing trajectory parameters that satisfy throughput constraints while minimizing scanner-induced defect costs.

%%% ----------------------------------------------------------------------
